% Generated draft: requires mathematical and source review.
\documentclass[11pt]{article}
\usepackage[margin=25mm]{geometry}
\usepackage{fontspec}
\setmainfont{Latin Modern Roman}
\usepackage{amsmath,amssymb,mathtools,unicode-math}
\setmathfont{Latin Modern Math}
\usepackage{enumitem,needspace,xurl,hyperref}
\setlength{\parindent}{0pt}
\setlength{\parskip}{5pt}
\setlength{\emergencystretch}{3em}
\newcommand{\N}{\mathbb{N}}
\newcommand{\Z}{\mathbb{Z}}
\newcommand{\Q}{\mathbb{Q}}
\newcommand{\R}{\mathbb{R}}
\newcommand{\C}{\mathbb{C}}
\newcommand{\F}{\mathbb{F}}
\newcommand{\A}{\mathbb{A}}
\newcommand{\PP}{\mathbb{P}}
\newcommand{\E}{\mathbb{E}}
\newcommand{\eps}{\varepsilon}
\newcommand{\dd}{\,\mathrm{d}}
\begin{document}
\section*{UnsolvedMath ID: 20001306}
\textbf{Source problem code:} catalogue entry
\subsection{Definitions and mathematical statement}
For a fixed dimension $n\geq3$, let $\mathcal K_s^n$ be the space of full-dimensional origin-symmetric convex bodies in $\R^n$, with Hausdorff distance. For $K\in\mathcal K_s^n$, its intersection body $IK$ has radial function
\[
 \rho_{IK}(u)=\operatorname{vol}_{n-1}(K\cap u^\perp),\qquad u\in S^{n-1}.
\]
For an origin-containing convex body $L$, its polar is $L^\circ=\{y:\langle x,y\rangle\leq1\ \forall x\in L\}$. A zonoid is a Hausdorff limit of finite Minkowski sums of line segments. Here $IK$ is convex in the origin-symmetric convex setting, so saying that $IK$ is the polar of a zonoid is equivalent to saying that $(IK)^\circ$ is a zonoid.

R. Schneider's requested Baire-category assertion is that
\[
 \{K\in\mathcal K_s^n:(IK)^\circ\text{ is not a zonoid}\}
 \quad\text{is residual in }\mathcal K_s^n,
\]
where residual means that its complement is meager (a countable union of nowhere dense sets). The source says ``most'' in this sense, not in a probability or volume sense. Its excerpt does not specify a dimension or topology; the fixed-$n\geq3$ Hausdorff formulation is the dimensionally qualified version supplied by the research aid. In dimension two the aid says the literal assertion is false: $IK$ is a rotated, scaled copy of $K$ and every origin-symmetric planar convex body is a zonoid.

The supplied aid reports that the exceptional set
\[
 E_n=\{K\in\mathcal K_s^n:(IK)^\circ\text{ is a zonoid}\}
\]
is closed. Thus the proposed generic assertion reduces to density of its open complement, equivalently to $E_n$ having empty interior. The aid also describes an affine-invariant version on the Banach--Mazur compactum of linear-equivalence classes; this does not replace the stated Hausdorff formulation with an unspecified new topology.

For completeness, the aid's dual obstruction can be stated with explicit integral normalizations. For $\varphi$ on $S^{n-1}$ define the unnormalized spherical Radon transform by
\[
 (R\varphi)(u)=\int_{S^{n-1}\cap u^\perp}\varphi(v)\,d\sigma_{n-2}(v).
\]
Then its notation is
\[
 f_K(u)=\frac{n-1}{R(\rho_K^{n-1})(u)}=h_{(IK)^\circ}(u),
\]
where $h_L(u)=\sup_{x\in L}\langle x,u\rangle$ is the support function. For a finite even signed Borel measure $\nu$, even meaning $\nu(A)=\nu(-A)$, write
\[
 (C\nu)(u)=\int_{S^{n-1}}|\langle u,v\rangle|\,d\nu(v).
\]
The reported certificate is
\[
 (IK)^\circ\text{ is not a zonoid}
 \quad\Longleftrightarrow\quad
 \exists\nu:\ C\nu\geq0\text{ pointwise},\quad \int_{S^{n-1}}f_K\,d\nu<0.
\]
The aid states that a strict witness persists under sufficiently small Hausdorff perturbations. These are attributed partial results; closedness and a certificate for a given example do not themselves establish the required density.
\subsection{Short English statement}
Schneider asks whether a typical symmetric convex body's intersection body fails to be the polar of a zonoid, with typical understood in the Baire-category sense. The supplied partial work makes the exceptional class closed and gives a strict separating-measure certificate, leaving density as the missing step. Dimension two is excluded because the literal assertion fails there.
\subsection{Sources}
\begin{itemize}
\item Local UnsolvedMath registry, permanent ID 20001306, exact source code catalogue entry; canonical \texttt{aim-convex-geometry-notes.json}, \texttt{notes[37]}, source item 0, tag ``problem'', workshop \emph{Fourier analytic methods in convex geometry}. Original attribution: R. Schneider. Supplied URL: \url{https://aimath.org/WWN/fourierconvex/fourierconvex.pdf}.
\item The original states the Baire-category meaning of ``most'' but supplies no separate bibliographic references. Its trailing isolated ``1'' is retained as extraction evidence rather than a mathematical condition. The verification consists only of ``The canonical record reads''; original remarks and literature field are empty.
\item Supplied research metadata: PARTIAL-PROGRESS, registry status ``partially\_solved'', statement status ``exact''; attempt 1, result ``partial\_result'', problem status at run ``open'', recorded rigor ``proved'', novelty confidence ``low'', candidate type ``dual obstruction lemma''. The research aid explicitly is machine generated, not peer reviewed, and requires independent verification before citation. No outside source was visited.
\end{itemize}
\end{document}
